\begin{topic}[id=tasmota_building,
             difficulty={Mittel},
             type={Systemanalyse + Praxisbezug},
             prerequisites={Grundlagen IoT und Gebäudeautomation; Basiswissen IP-Netze hilfreich},
             practice={gut möglich},
             owner={Dozententeam},
             updated={2026-04-13},
             tags={ Matter, Thread, Gebäudeautomation, Interoperabilität, Provisioning, Security, Home Assistant }]{Matter und Thread in der Gebäudeautomation}

\TopicDescription{Für die Gebäudeautomation wird 2026 weniger interessant, wie einzelne Open-Source-Firmwares konfiguriert werden, sondern wie interoperable, lokale und zukunftsfeste Integrationspfade aussehen. Das Thema fokussiert daher Matter als Anwendungsstandard und Thread als IP-basiertes Mesh-Netz für vernetzte Geräte. Betrachtet werden Provisioning und Commissioning, Fabric- und Border-Router-Konzepte, lokale Steuerung ohne Cloud-Zwang, Security-Grundlagen sowie die praktische Einordnung gegenüber bestehenden MQTT-, Home-Assistant- und Bestandssystemen. Tasmota darf höchstens als Brownfield- oder Open-Source-Kontext vorkommen; der Kern ist die technische Bewertung von Interoperabilität, Betriebsmodell und Integrationsgrenzen.}

\TopicLeitfragen{\item Welche Probleme lösen Matter und Thread gegenüber klassischen Smart-Home-Insellösungen wirklich?
\item Wie wirken sich Commissioning, lokale IP-Integration und Security-Mechanismen auf Betrieb und Wartbarkeit aus?
\item Wie integriert man Matter/Thread sinnvoll in bestehende MQTT- oder Home-Assistant-Landschaften, ohne neue Komplexität zu verstecken?
}
\TopicScope{\item Kein Reverse-Engineering proprietärer Geräte.
\item Keine Elektroinstallation.
}
\TopicPractice{Optionale Praxisidee: Architekturvergleich für ein kleines Raum-Szenario mit Matter/Thread und einer Bestandsintegration über Home Assistant oder Bridge-Konzept.}

\TopicKeywords{Matter, Thread, Gebäudeautomation, Interoperabilität, Provisioning, Security, Home Assistant}

\nocite{matterGoogleHome,openThreadPrimer,homeAssistant}
\end{topic}
